\documentclass[11pt,a4paper]{article}

\input{../preambule}

\begin{document}

\title{Semaine 1}
\author{}
\date{21/09 - 27/09 }

\maketitle


\begin{exercise}
Soit $n\in\N^*$, trouver les couples $(k,l)\in[\![0;n]\!]$ tels que $k\lt l$ et $\binom{n}{k}=\binom{n}{l}$
\end{exercise}

\begin{exercise}
Soit $A\subset\R^+$ non majorée. On pose $B=\set{\frac{a}{n},a\in A, n\in\N^*}$.\\
Montrer que $B$ est dense dans $\R^+$
\end{exercise}

\begin{exercise}
Soit $0\lt\varepsilon\lt1, n\in\N^*, x_1,\ldots,x_n\in\R^+$ tels que $\forall (i,j)\in\intset{1,n}^2, x_ix_j\leq\varepsilon^{\abs{i-j}}$.\\
Montrer que $\displaystyle\sum_{k=1}^nx_k\leq \frac{1}{1-\sqrt{\varepsilon}}$
\end{exercise}

\begin{exercise}
Soit $n\in\N^*$.
Soient $p_1,\ldots,p_n\in\R_+$ tels que
$\displaystyle\sum_{k=1}^np_k=1$
et $\omega_1,\ldots,\omega_n\in\R$. \\
Soit f la fonction définie sur $\R$ par
$f(x)=\displaystyle\sum_{k=1}^np_k\cos{(\omega_kx)}$.\\
Démontrer que
$f^2(x)\le\frac{1+f(2x)}{2}$.
\end{exercise}

\begin{exercise}[Inégalité de Weierstrass]
Soit $n>0$ un entier et, $u_1,\ldots,u_n$ et $v_1,\ldots,v_n$ deux familles de complexes tels que $|u_k|,|v_k|\le1$ pour tout entier $1\le k\le n$. Démontrer que pour tout $n\in\N^*$,
\[\left|\prod_{k=1}^nu_k-\prod_{k=1}^nv_k\right|\le\sum_{k=1}^n|u_k-v_k|.\]
\end{exercise}

\begin{exercise}
Soit $n\in\N^*$, calculer $\displaystyle\sum_{k=0}^n\binom{n}{k}^2$.
\end{exercise}

\begin{exercise}
Soient $5$ points sur la surface d'une sphère. Montrer qu'il existe un hémisphère contenant au moins $4$ de ces points.
\end{exercise}

\begin{exercise}
Montrer que :
$$\forall n\in\N^*, \sum_{k=1}^nk^5+\sum_{k=1}^nk^7=2\left(\sum_{k=1}^nk\right)^4$$
\end{exercise}

\begin{exercise}
\begin{questions}
\item Montrer que pour tout $n\in\N^*, \floor{\sqrt{n^4+2n^3+3n^2+1}}=n^2+n$.
\item Déterminer l'ensemble des entiers $n\in\N^*$ tels que $n^4+2n^3+3n^2+1$ est un carré parfait.
\end{questions}
\end{exercise}

\begin{exercise}
\begin{questions}
\item Soit $A\subset\R$ non vide.
\begin{questions}
\item Justifier pour $x\in\R$ l'existence du réel $d(x,A)=\inf_{a\in A}|x-a|$.
\item Calculer $d(x,A)$ pour tout $x\in A$.
\item Montrer que pour tous $x,y\in\R$, on a $|d(x,A)-d(y,A)|\le|x-y|$.
\end{questions}
\item On suppose ici que $A=\Q\cap[0,1[$. Déterminer $d(x,A)$ pour tout $x\in\R$ et tracer le graphe de la fonction $x\mapsto d(x,A)$.
\end{questions}
\end{exercise}

\begin{exercise}
Soit $n \ge 2$ un entier et $a \in ]0,1[$. Montrer que $(1-a)^n \le \frac{1}{1+na}$.
\end{exercise}

\begin{exercise}
Soit $n\in\N^*$ et $\theta\in\R$. Calculer $\displaystyle\sum_{k=0}^{n-1}\cos^3(k\theta)$.
\end{exercise}

\begin{exercise}
Soit $z\in \mathbb{U} - \{1\}$. Montrer qu'il existe $n\in\N$ tel que $\left| z^n - 1 \right| \ge \sqrt{3}$.
\end{exercise}

\begin{exercise}
Soit $n\in\N$. Calculer $\displaystyle\sum_{k=0}^{n}\binom{n-k}{k}$.
\end{exercise}

\begin{exercise}
Soit $z\in \C$. Notons $\theta$ un argument de $z$ dans $]-\pi; \pi]$.
Montrer que $|z-1|\le ||z|-1|+|z\theta|$.
\end{exercise}

\begin{exercise}
Soit $n\in\N^*$.
\begin{questions}
\item Montrer qu'il existe d'uniques $x_n$ et $y_n$ entiers strictement positifs tels que $(3+2\sqrt{2})^n=x_n+\sqrt{2}y_n$.
\item Exprimer $x_{n+1}$ et $y_{n+1}$ à partir de $x_{n}$ et $y_{n}$. En déduire que les deux suites sont strictement croissantes.
\item Montrer que $(3-2\sqrt{2})^n=x_n-\sqrt{2}y_n$.
\item En déduire que l'équation $x^2-2y^2=1$ admet une infinité de solutions entières.
\end{questions}
\end{exercise}

\begin{exercise}
\begin{questions}
\item Soit $(x_n)_{\in\N}$ une suite de rééls telle que : $$\forall n \ge0, \sum_{k=0}^nx_k^3=\left(\sum_{k=0}^nx_k\right)^2$$
Montrer que pour tout $n\ge0$, il existe un entier naturel $m$ tel que : $$\sum_{k=0}^nx_k=\frac{m(m+1)}{2}$$
\item On pose $S_{n,p}=\displaystyle\sum_{k=0}^nk^p$. Déterminer les entiers naturels non nuls $p$ tels que, pour tout $n\ge0$, $S_{n,p}$ soit le carré d'un entier naturel.
\end{questions}
\end{exercise}

\begin{exercise}
Calculer :
$$
\sum_{1 \leq p, q \leq n} \min(p,q)
$$
\end{exercise}

\begin{exercise}
Déterminer tous les $(x_1, \dots, x_n) \in \mathbb R^n$ tels que :
$$
\sum_{k=1}^n x_k^2 = \sum_{k=1}^n x_k = n
$$
\end{exercise}

\begin{exercise}
Soit $n \geq 2$ et $a,b \in \C$. Montrer que :
$$
\abs a + \abs b \leq \frac 2 n \sum_{\omega \in \U_n} \abs{a + \omega b}
$$
\end{exercise}

\begin{exercise}
On définit, pour toute partie $A$ de $\N$, la densité de $A$ valant $d(A)=\inf_{a\in A}{\frac{|A\cap [1;n]|}{n}}$.
1) Soit $A$ une partie de $\N$ contenant $0$ de densité au moins $0,5$. Montrer que tout naturel s'écrit comme somme de deux éléments de $A$.
2) $0\in A$ est-il une condition nécessaire, et "à quel point" ?
3) Soient $A,B \subset N$ contenant toutes deux $0$. On pose $A+B={a+b, a\in A, b\in B}$. Montrer que $1-d(A+B)\le (1-d(A))(1-d(B))$.

\end{exercise}








\end{document}